\documentclass[11pt]{article}
\usepackage[T1]{fontenc}
\usepackage[utf8]{inputenc}
\usepackage{lmodern}
\usepackage[margin=1in]{geometry}
\usepackage{amsmath,amssymb,amsthm,mathtools,microtype,booktabs,xcolor}
\usepackage[colorlinks=true,linkcolor=blue!45!black,citecolor=blue!45!black,urlcolor=blue!45!black]{hyperref}
\newtheorem{theorem}{Theorem}[section]
\newtheorem{lemma}[theorem]{Lemma}
\newtheorem{proposition}[theorem]{Proposition}
\newtheorem{corollary}[theorem]{Corollary}
\theoremstyle{remark}\newtheorem{remark}[theorem]{Remark}
\newcommand{\R}{\mathbb R}\newcommand{\E}{\mathbb E}\newcommand{\Prb}{\mathbb P}
\newcommand{\relu}{\operatorname{ReLU}}\newcommand{\gelu}{\operatorname{GELU}}
\newcommand{\norm}[1]{\left\lVert#1\right\rVert}
\newcommand{\cX}{\mathcal X}\newcommand{\cF}{\mathcal F}\newcommand{\cG}{\mathcal G}
\newcommand{\err}{\operatorname{err}}\newcommand{\Rad}{\widehat{\mathfrak R}}
\setlength{\emergencystretch}{2em}
\hypersetup{pdftitle={Width-independent learning bounds for ReLU, GELU, and other Lipschitz activations},pdfauthor={ChatGPT-6 Astra; Claude Fable 5 and Opus 5; Kaarel Hanni; Dmitry Vaintrob}}
\title{Width-independent learning bounds for ReLU, GELU,\protect\\ and other Lipschitz activations}
\author{ChatGPT-6 Astra \and Claude Fable 5 and Opus 5 \and Kaarel H\"anni \and Dmitry Vaintrob\thanks{Authors are listed alphabetically.}}
\date{October 22, 2026}
\begin{document}
\maketitle
\begin{abstract}
We prove polynomial sample-complexity bounds at fixed depth for Bayesian neural networks with any fixed globally Lipschitz activation that vanishes at zero, including ReLU, exact GELU, SiLU, and tanh. The target has a width-$C$ implementation with weights bounded by $C$ and strict signed output margin greater than $1/2$ on the Boolean cube. The learner's Gaussian prior is conditioned on the same absolute margin at every input, then on agreement with noiseless training labels. One polynomial sample bound controls posterior average sign error simultaneously at every learner width $N\ge C$ and every mean-field-or-stronger prior scale. Here the scale parameter $\sigma^2$ is the internal weight variance, and the covered range is $0<\sigma^2\le1/N$. This includes all power scalings $\sigma^2=N^{-\gamma}$ with $\gamma\ge1$, without deterioration as $\gamma$ approaches one. The proof places all but a small relative amount of prior mass in a finite class of Boolean functions. Spectral truncation supplies that class at small scales; at mean field, a bounded Gaussian part and controlled tail atoms yield finite covariance reconstructions. The counting class allows neuron-specific Lipschitz activations, while the Bayesian learner uses its one fixed activation. We also give a simpler strong-prior proof and an exact same-depth GELU implementation of layered Boolean circuits with at most twice the width. The sample polynomials have poor depth dependence.
\end{abstract}

\section{Introduction}
We study noiseless Boolean learning when a small neural network implements the target and the Bayesian learner may be much wider. Both networks have the same depth and activation. Before seeing data, the learner retains only parameters whose outputs stay strictly away from zero at every Boolean input; training then retains the functions that agree with all observed labels. The question is whether the number of examples needed to control this posterior can depend on the teacher's width and weight bound while remaining independent of the learner's width.

Bounded-activation proofs can compress hidden states by sampling neurons. ReLU and GELU require additional control because individual hidden states can be arbitrarily large, and GELU does not permit the exact rescaling identities available for ReLU.

We prove that neither boundedness nor homogeneity is necessary for a fixed-depth result. A globally Lipschitz activation $\phi$ satisfying $\phi(0)=0$ is enough. The proof controls the product of layer Frobenius norms, embeds the resulting networks into an auxiliary class with neuron-specific Lipschitz activations, and bounds that larger class using a one-dimensional activation cover. This gives a finite class of globally margin-separated Boolean functions at each fixed energy, simultaneously across all widths.

That capacity statement has two applications. In the strong-prior limit, taken at each fixed learner width, Gaussian rates suppress functions whose minimum energy exceeds the target's. At every mean-field-or-stronger scale, a separate argument approximates margin-separated draws outside a small exceptional event by auxiliary networks of controlled complexity. For small scales, spectral truncation discards matrix components with small operator norm. At larger scales, including mean field, bounded exponential tilts describe the regular part of the hidden covariance, while a reproducing-kernel argument retains the contribution of rare large neurons. Comparing the exceptional mass with the target mass then gives one polynomial sample bound across the entire range of scales.

These are information-theoretic guarantees for exact Bayesian conditioning, not claims about efficient posterior computation. They concern iid noiseless learning on the Boolean cube and the specified global-margin conditioning. The polynomial degree depends on depth; a polynomial-depth theorem requires additional ideas or restrictions. Established tools used below include Rademacher contraction and peeling, Gaussian concentration, and finite-class Occam learning~\cite{golowich,blumer,borell}. We give the particular estimates and normalization needed here.

\section{Model and the results}
Fix input dimension $d\ge1$, depth $L\ge2$, and $H=L-1$ hidden layers. Let $\cX=\{-1,1\}^d$, $\bar x=(1,x_1,\ldots,x_d)$, and fix an activation $\phi:\R\to\R$ with $\phi(0)=0$ and Lipschitz constant at most $\lambda\ge1$. A width-$N$ network is
\[
 h^0(x)=\bar x,\qquad h^s(x)=\phi(W^s h^{s-1}(x))\quad(1\le s<L),
 \qquad z_w(x)=v^\top h^{L-1}(x),
\]
where activation is coordinatewise and every hidden layer has width $N$. The appended input coordinate supplies first-layer offsets. There are no separate internal biases: any constant wire used to supply later offsets must be implemented by the network and counted in its width. Define the weighted Gaussian energy
\begin{equation}\label{eq:energy}
 \kappa_N(w)=\frac dN\norm{W^1}_F^2+
 \sum_{s=2}^{L-1}\norm{W^s}_F^2+N\norm v_2^2.
\end{equation}
The raw prior $\mu_{N,\sigma}$ has independent Gaussian coordinates with variances
\[
 (N/d)\sigma^2\quad\text{in the first layer},\qquad
 \sigma^2\quad\text{internally},\qquad
 \sigma^2/N\quad\text{in the readout}.
\]
Its density is proportional to $\exp[-\kappa_N(w)/(2\sigma^2)]$. Prescribed exponent $\gamma$ means $\sigma^2=N^{-\gamma}$.

For a Boolean function $g$, set
\[
 U_{N,g}=\{w:g(x)z_w(x)>1/2\text{ for every }x\in\cX\},
 \qquad U_N=\bigcup_g U_{N,g}.
\]
Condition the raw prior on $U_N$ before seeing data. This induces the function prior
\[
 \pi_{N,\sigma}(g)=\frac{\mu_{N,\sigma}(U_{N,g})}{\mu_{N,\sigma}(U_N)}.
\]
Given iid examples $S_n=((X_i,f(X_i)))_{i=1}^n$ with $X_i\sim D$, condition again on agreement with all labels, obtaining $\Pi_{N,\sigma}(\cdot\mid S_n)$. Its average sign error is
\[
 R_D(\Pi)=\E_{g\sim\Pi}\Prb_{X\sim D}\{g(X)\ne f(X)\}.
\]
The distribution $D$ may be any distribution on the cube.

Assume $f$ has a width-$C$, depth-$L$ teacher with all individual weights bounded in absolute value by $C$, satisfying $f(x)z_*(x)>1/2$ for every input. This is the only teacher assumption. For an arbitrary activation it need not hold for every Boolean circuit; we verify it with a small implementation for ReLU and exact GELU in Section~\ref{sec:gelu}.

\begin{theorem}[Summary of learning guarantees]\label{thm:summary}
For fixed $L$ and $\lambda$, there is a sample bound polynomial in
$d,C,\varepsilon^{-1},\log(1/\delta)$, independent of learner width,
that gives $R_D(\Pi_{N,\sigma}(\cdot\mid S_n))\le\varepsilon$ with
probability at least $1-\delta$, simultaneously for every $N\ge C$ and
$0<\sigma^2\le1/N$. Thus one polynomial works for every prescribed
exponent $\gamma\ge1$ in $\sigma^2=N^{-\gamma}$, and also covers the
strong-prior limit at every fixed width.

Theorem~\ref{thm:uniform-all-scales} gives the quantitative bound through
constructive polynomial budgets. Theorem~\ref{thm:strong} gives a simpler
strong-prior argument with a smaller explicit bound.
\end{theorem}

The probability statements are over the training inputs. The posterior may mix tied implementations or tied minimum-energy functions. No uniqueness is assumed. If $p_S(x)$ is the posterior probability of the correct label, then
\begin{equation}\label{eq:confidence}
 D\{x:p_S(x)\le q\}\le \frac{R_D(\Pi)}{1-q}\qquad(0<q<1).
\end{equation}
Thus using $\varepsilon/3$ in an average-risk bound gives correct-label probability strictly greater than $2/3$ except on an input set of measure at most $\varepsilon$. Posterior-majority error is at most twice the average sign error.

Examples include ReLU and tanh with $\lambda=1$; exact GELU $\phi(t)=t\Phi(t)$, where $\Phi$ is the standard normal distribution function, and SiLU $\phi(t)=t/(1+e^{-t})$ with $\lambda=2$; and centered softplus $\phi(t)=\log(1+e^t)-\log2$ with $\lambda=1$. A nonzero activation offset is not covered by the theorem as stated.

\section{A capacity estimate without boundedness or homogeneity}
The next estimates turn an energy bound into a bound on the number of Boolean functions that can be represented with a global margin. To make normalization possible without assuming homogeneity, we count a larger class in which each neuron can have its own Lipschitz activation. This enlargement is used only in the proof.
Let $\Psi$ be the family of all globally $1$-Lipschitz real functions $\psi$ with $\psi(0)=0$.

\begin{lemma}[A scalar activation cover]\label{lem:activation-net}
For $R>0$ and $0<\eta\le R$, the restrictions of $\Psi$ to $[-R,R]$ have a supremum-norm $\eta$-net with $M_\eta$ centers that themselves belong to $\Psi$, satisfying
\[
 \log(2M_\eta)\le14R/\eta.
\]
\end{lemma}
\begin{proof}
Put $J=\lceil3R/\eta\rceil$ and $\Delta=R/J$. Round the values of $\psi$ at the grid points $j\Delta$, $-J\le j\le J$, to the nearest multiple of $\Delta$. The value at zero has code zero, and successive rounded integers differ by at most two. There are at most $5^{2J}$ codes. From each nonempty code choose one actual member of $\Psi$ as its representative. Functions with the same code differ by at most $\Delta$ at grid points and at most $3\Delta\le\eta$ everywhere on the interval. Finally
\[
 \log(2M_\eta)\le\log2+2(3R/\eta+1)\log5\le14R/\eta.
\]
Choosing actual representatives preserves the Lipschitz constant exactly.
\end{proof}

\begin{lemma}[Convex contraction]\label{lem:contraction}
For a bounded set $T\subset\R^m$, a fixed $\psi\in\Psi$, and $a>0$,
\[
 \E_\epsilon\exp\left(a\sup_{t\in T}\sum_i\epsilon_i\psi(t_i)\right)
 \le \E_\epsilon\exp\left(a\sup_{t\in T}\sum_i\epsilon_i t_i\right).
\]
The same holds with any convex nondecreasing function in place of the exponential.
\end{lemma}
\begin{proof}
Replace one Rademacher coordinate at a time. With the other coordinates fixed, write $a_t$ for their contribution and
\[
 X_\pm=\sup_t(a_t\pm\psi(t_k)),\qquad Y_\pm=\sup_t(a_t\pm t_k).
\]
Because $|\psi(t_k)|\le|t_k|$, one has $\max(X_+,X_-)\le\max(Y_+,Y_-)$. Also
\[
 X_++X_-\le\sup_{t,u}(a_t+a_u+|t_k-u_k|)=Y_++Y_-.
\]
For two numbers, these maximum and sum inequalities imply the corresponding inequality for the sum of any convex nondecreasing function. Average over the coordinate and iterate. Suprema follow by finite approximation.
\end{proof}

Let $\cG_s$ consist of all vector-valued feature maps of depth $s$ whose layer matrices have Frobenius norm at most one and whose neurons use arbitrary, independently chosen members of $\Psi$. Widths can be any finite integers. Put $\cG_0=\{x\mapsto x\}$. All feature vectors have norm at most $R$ on inputs of norm at most $R$.

\begin{lemma}[Width-independent feature complexity]\label{lem:feature-rad}
For $H\ge1$, inputs $x_1,\ldots,x_m$ of Euclidean norm at most $R$, and independent uniform signs $\epsilon_i$,
\[
 \frac1m\E_\epsilon\sup_{h\in\cG_H}
 \left\|\sum_i\epsilon_i h(x_i)\right\|_2
 \le8HRm^{-1/3}.
\]
\end{lemma}
\begin{proof}
Write $Z_s=\sup_{h\in\cG_s}\|\sum_i\epsilon_i h(x_i)\|_2$. Decompose each row of the final matrix as $r_j u_j$, where $r_j\ge0$, $\|u_j\|_2\le1$, and $\sum_jr_j^2\le1$. For $r_j>0$, the map $t\mapsto\psi_j(r_jt)/r_j$ is again in $\Psi$; zero rows contribute zero. Consequently
\[
 Z_s\le\sup_{h\in\cG_{s-1},\,\|u\|_2\le1,\,\psi\in\Psi}
 \left|\sum_i\epsilon_i\psi(u\cdot h(x_i))\right|.
\]
Replace $\psi$ by a center from Lemma~\ref{lem:activation-net}, at an additive cost $m\eta$. For each center, handle the absolute value by its two signs and apply Lemma~\ref{lem:contraction}. Since $\sup_{h,u}\sum_i\epsilon_i u\cdot h(x_i)=Z_{s-1}$, taking a finite union gives
\[
 \E e^{aZ_s}\le2M_\eta e^{am\eta}\E e^{aZ_{s-1}}\qquad(a>0).
\]
At depth zero, $\E Z_0\le R\sqrt m$. Changing one sign changes $Z_0$ by at most $2R$, so the bounded-differences exponential estimate gives
\[
 \E e^{aZ_0}\le\exp(aR\sqrt m+a^2mR^2/2).
\]
Iterate $H$ times and optimize $a$ to obtain
\[
 \E Z_H\le Hm\eta+R\sqrt m+R\sqrt{2mH\log(2M_\eta)}.
\]
Taking $\eta=Rm^{-1/3}$ and using $1+H+\sqrt{28H}\le8H$ proves the claim. No bound here depends on the number of neurons. All suprema are well-defined on the finite Rademacher sample space; equivalently one can take finite subclasses and pass to monotone limits.
\end{proof}

\begin{proposition}[Product-norm bound for Lipschitz activations]\label{prop:product-rad}
Consider networks with $H=L-1$ hidden layers, with a possibly different zero-at-zero $\lambda$-Lipschitz activation at each neuron, and with product of layer Frobenius norms at most $P$. Their empirical Rademacher complexity is at most
\[
 8H\lambda^H PRm^{-1/3}
\]
on any $m$ inputs of norm at most $R$, independently of the finite widths.
\end{proposition}
\begin{proof}
Write $m_s=\|W^s\|_F$, treating the readout as layer $L$, and let $\phi_{s,j}$ be the activation of neuron $j$ in hidden layer $s$. If some $m_s=0$, the network is zero. Otherwise put $q_0=1$, $q_s=\lambda^s\prod_{j\le s}m_j$, and $g_s=h_s/q_s$. With $A_s=W^s/m_s$ one has the exact coordinate identities
\[
 g_{s,j}=\psi_{s,j}((A_sg_{s-1})_j),\qquad
 \psi_{s,j}(t)=\frac{\phi_{s,j}(m_sq_{s-1}t)}{\lambda m_sq_{s-1}}\in\Psi.
\]
The output is $\lambda^H\prod_sm_s$ times a norm-one linear readout of $g_H$. Apply Lemma~\ref{lem:feature-rad}.
\end{proof}

\begin{remark}
This normalization embeds a fixed-activation network into a larger auxiliary class by changing the scalar activations. It therefore applies to GELU without a homogeneity assumption. The prior and teacher continue to use the single fixed activation $\phi$; only the auxiliary representation uses the normalized functions.
\end{remark}

For $K>0$ let $\cF(K)$ contain every Boolean function admitting a depth-$L$, globally strict-margin-$1/2$ implementation of energy at most $K$, at some finite equal hidden width, allowing each hidden neuron its own activation vanishing at zero with Lipschitz constant at most $\lambda$. This is an auxiliary counting class, larger than the fixed-$\phi$ learner's function class. Proposition~\ref{prop:product-rad} and its proof apply unchanged to this neuron-specific class. Define
\begin{equation}\label{eq:capacity-budget}
 v(K)=\left\lceil\left[16H\lambda^H\sqrt{\frac{d+1}{d}}
                   \left(\frac KL\right)^{L/2}\right]^3\right\rceil,
 \qquad \mathcal B(K)=v(K)\log(2^d+1).
\end{equation}

\begin{corollary}[Counting globally margin-separated functions]\label{cor:count}
The common class across all widths satisfies $\log|\cF(K)|\le\mathcal B(K)$.
\end{corollary}
\begin{proof}
The squared norms of $\sqrt{d/N}W^1,W^2,\ldots,W^{L-1},\sqrt N v^\top$ sum to $\kappa_N$ and their product equals $d\prod_s\|W^s\|_F^2$. Arithmetic--geometric mean yields
\[
 \prod_s\|W^s\|_F\le d^{-1/2}(K/L)^{L/2}.
\]
The appended inputs have norm $\sqrt{d+1}$. Shattering $m$ cube inputs with global margin $>1/2$ forces the empirical Rademacher complexity of the underlying real-output class to be at least $1/2$. Proposition~\ref{prop:product-rad} bounds its VC dimension by $v(K)$. The Sauer bound on the $2^d$-point domain gives $|\cF(K)|\le(2^d+1)^{v(K)}$.
\end{proof}

\section{The strong-prior learning theorem}
The teacher gives a width-independent energy budget. At each fixed learner width, sufficiently small Gaussian variance makes functions above that budget negligible relative to the target. The common finite function class from the preceding section then supplies the sample bound.
\begin{lemma}[Exact width enlargement]\label{lem:clone}
The target teacher can be implemented at every width $N\ge C$ with the same real-valued output and energy at most
\[
 B_*=d(d+1)C^2+LC^4.
\]
\end{lemma}
\begin{proof}
Let $k=\lfloor N/C\rfloor$. Make $k$ copies of each hidden neuron, copy the first-layer incoming weights, divide all internal weights by $k$ between each pair of copied populations, and divide the readout weights by $k$. Fill remaining slots with zero neurons; $\phi(0)=0$ makes padding harmless. Each copy computes exactly the original neuron's activation. The first-layer energy is at most $d(d+1)C^2$, internal Frobenius energies are unchanged and at most $C^4$ each, and readout energy is at most $2C^4$ because $N/k<2C$. This gives the claimed sum.
\end{proof}

\begin{lemma}[Finite-dimensional Gaussian rate]\label{lem:rate}
For fixed $N$ and nonempty $U_{N,g}$, put $k_N(g)=\inf_{w\in U_{N,g}}\kappa_N(w)$. Then
\[
 \lim_{\sigma\downarrow0}2\sigma^2\log\mu_{N,\sigma}(U_{N,g})=-k_N(g).
\]
\end{lemma}
\begin{proof}
In weighted coordinates the raw law is an isotropic Gaussian. A small ball around any point of the open event gives the lower bound with rate equal to that point's squared norm, and then to the infimum. The closed-set Gaussian upper bound applies to its closure. A point of this closure has signed outputs at least $1/2$ everywhere; multiplying its readout by any factor greater than one places it in $U_{N,g}$ with arbitrarily small increase in energy. Hence the closure and the open event have equal energy infima. The elementary Gaussian upper bound can equivalently be obtained from the chi-square tail outside the ball with squared radius $k_N(g)$.
\end{proof}

\begin{theorem}[Strong-prior learning]\label{thm:strong}
Let $B_*$ be as in Lemma~\ref{lem:clone}. For every $N\ge C$ and $0<\varepsilon<1$, there is $\sigma_0(N)>0$ such that for $0<\sigma<\sigma_0(N)$,
\[
 n\ge\frac2\varepsilon\left[\mathcal B(B_*+1)+\log(1/\delta)\right]
\]
implies $R_D(\Pi_{N,\sigma}(\cdot\mid S_n))\le\varepsilon$ with sample probability at least $1-\delta$, for every $D$ and $0<\delta<1$. The threshold can depend on $d,L,\lambda,N$, the teacher, and $\varepsilon$, but not on $D,n,\delta$. A common sample event works at every $N\ge C$, using the width-specific variance thresholds.
\end{theorem}
\begin{proof}
At fixed width, retain functions satisfying $k_N(g)\le k_N(f)$. They are all in $\cF(B_*+1)$ by Lemma~\ref{lem:clone} and the definition of infimum. There are finitely many Boolean functions, so Lemma~\ref{lem:rate} makes the total raw mass of the remaining realizable functions divided by the target mass tend to zero. Choose $\sigma_0(N)$ so the ratio is at most $\varepsilon/2$. Global-margin normalization cancels. On every noiseless training history, the posterior denominator includes the target; thus posterior mass outside the retained class remains at most $\varepsilon/2$.

By a finite-class union bound, with probability at least $1-\delta$, every member of $\cF(B_*+1)$ that agrees with all $n$ labels has true error at most $\varepsilon/2$. The stated sample bound suffices since a function of error greater than $\varepsilon/2$ survives with probability at most $e^{-n\varepsilon/2}$. On this common event, posterior average error is at most $\varepsilon$. This event does not depend on width.
\end{proof}

At fixed $L,\lambda$, this bound is polynomial in the indicated variables; its crude power of $C$ is at most $6L$. Positive homogeneity allows a sharper estimate for ReLU, but is unnecessary for the existence of a polynomial. The strong-prior limit is taken at fixed width before invoking the conclusion, and does not assert that one variance threshold works at all widths.

% Standalone input fragment. Expects amsmath,amssymb,amsthm, proposition,
% lemma,theorem environments. All macros below use standard LaTeX.
\section{Target mass and spectral compression}
\label{sec:allgamma-spectral}
This section extends the Gaussian spectral argument to unbounded
activations. Let $\phi(0)=0$ and let $\phi$ be $\lambda$-Lipschitz; replace
$\lambda$ by $\max\{1,\lambda\}$, so throughout $\lambda\ge1$.
For $\gamma\ge1$, the independent coordinate variances are
$N^{1-\gamma}/d$ in the first layer, $N^{-\gamma}$ internally, and
$N^{-1-\gamma}$ at the readout. Write $U_{N,g}$ for the globally
strict-margin event $g(x)z_w(x)>1/2$ for every cube input.
The normalized energy is
\[
 \kappa_N(w)=\frac dN\|W^1\|_F^2+
 \sum_{s=2}^{L-1}\|W^s\|_F^2+N\|v\|_2^2.
\]
No uniform bound on $\phi$ or on individual hidden activations is imposed.

\subsection{A target cylinder for linearly growing activations}
Let a depth-$L$, width-$C$ teacher have every weight bounded by $C$ and
strict output margin $1/2$. Multiplying only its readout by three gives
strict margin $3/2$. Set
\begin{align*}
 M&=(d+1)(\lambda C^2)^{L-1},&
 b&=2+\sqrt{2\log(2L)},\\
 B&=2+\lambda(\sqrt2 C^2+\sqrt C+b),&
 E&=\lambda M(1+LC)B^{L-2},\\
 \rho&=\left[4\{(3\sqrt2 C^2+1)E+M\}\right]^{-1},\\
 B_0&=d(d+1)C^2+2(L-2)C^4+36C^4,&
 D_0&=d+2+(L-2)C,\\
 A&=2B_0+4D_0/\rho+2.
\end{align*}
These quantities are independent of learner width and polynomially bounded
in $d,C$ for fixed $L,\lambda$. In particular $0<\rho\le1$ and $A\ge1$.

\begin{proposition}[Unbounded-activation target mass]
\label{prop:unbounded-target-mass}
At every $N\ge C$ and every $\gamma\ge1$, the raw Gaussian probability
of the target event satisfies
\[
 \mu_{N,\gamma}(U_{N,f})\ge \exp(-AN^\gamma/2).
\]
\end{proposition}
\begin{proof}
Partition the $N$ neurons into $C$ groups of sizes $k_a$ differing by at
most one. Then $N/k_a\le2C$ and $k_a/k_b\le2$. Clone the teacher by
copying each first-layer row throughout its group, replacing an internal
weight $B^s_{ba}$ by $B^s_{ba}/k_a$ on edges from group $a$ to group $b$,
and replacing readout weight $v_b$ by $v_b/k_b$. This preserves the exact
real-valued function for any activation. Every cloned hidden state is
group-constant. Let $P$ project onto that subspace.

Because $|\phi(u)|\le\lambda|u|$, all teacher hidden coordinates are
bounded in magnitude by $M$: the first layer is bounded by
$\lambda C(d+1)$, and each next hidden layer multiplies this bound by at
most $\lambda C^2$. The cloned internal operator norms are at most
$\sqrt2 C^2$, and the scaled readout norm is at most
$\sqrt N\|v_*\|_2\le3\sqrt2 C^2$.

Put $t=N^{(1-\gamma)/2}\le1$. In the first layer constrain each
$\sqrt d W^1_{ij}$ within $\rho t$ of its cloned value. In each internal
row, constrain its $C$ coordinates along the normalized group indicators,
after multiplying the row by $\sqrt N$, within $\rho t$ of their cloned
values. Finally constrain each $Nv_j$ within $\rho t$ of its cloned
value. All these are independent scalar Gaussian coordinates of variance
$t^2$. Write the resulting internal matrices as
\[
 W^s=W_*^s+R^s+E^s,\qquad R^s=R^sP,\quad E^sP=0.
\]
The constrained error obeys $\|R^s\|_F\le\rho t\sqrt C$; the free
part is $(t/\sqrt N)G^s(I-P)$, independently of the constrained
coordinates, with standard Gaussian $G^s$. The Gaussian operator-norm
bound and a union bound give conditional probability at least $1/2$ that
$\|E^s\|_{\rm op}\le bt$ for all internal layers.

On this event put
$e_s(x)=N^{-1/2}\|h_w^s(x)-h_*^s(x)\|_2$. Uniformly over the cube,
\[
 e_1\le\lambda\rho(d+1),\qquad
 e_s\le B e_{s-1}+\lambda\rho\sqrt C\,M.
\]
The second estimate uses $E^sh_*^{s-1}=0$, the bound on the teacher hidden
states, and Lipschitz continuity. It does not bound the perturbed states
by one. Since $M\ge d+1$ and $B\ge1$, induction gives
$e_{L-1}\le\rho E$. The readout error has scaled norm at most $\rho$,
so
\[
 \sup_x|z_w(x)-z_*(x)|
 \le (3\sqrt2 C^2+1)e_{L-1}+\rho M\le1/4.
\]
Thus every such draw lies in $U_{N,f}$ with room to spare.

There are $ND_0$ constrained coordinates and their squared centers sum
to at most $NB_0$. For $Z\sim N(0,t^2)$,
\[
 \mathbb P\{|Z-a|\le\rho t\}
 \ge\frac{2\rho}{\sqrt{2\pi}}e^{-a^2/t^2-\rho^2}.
\]
Multiplying these bounds, retaining the independent operator-norm event,
and using $N\le N^\gamma$ gives
\[
 -\log\mu_{N,\gamma}(U_{N,f})
 \le N^\gamma B_0+ND_0\log\frac{e\sqrt{2\pi}}{2\rho}+\log2
 \le AN^\gamma/2.
\]
Here $\log(e\sqrt{2\pi}/(2\rho))\le2/\rho$ for $0<\rho\le1$.
\end{proof}

\subsection{Spectral truncation with controlled hidden RMS norms}
\begin{lemma}[Gaussian spectral tail]
\label{lem:unbounded-spectral-tail}
If $W$ is an $N\times N$ matrix with independent $N(0,\sigma^2)$ entries,
$a\ge8\sigma\sqrt N$, and $W_{>a}$ retains singular values greater than
$a$, then
\[
 \mathbb P\{\|W_{>a}\|_F^2>K\}\le e^{-K/(8\sigma^2)}.
\]
\end{lemma}
\begin{proof}
The Gaussian operator-norm inequality
$\mathbb P\{\|W\|_{\rm op}>2\sigma\sqrt N+\sigma u\}\le e^{-u^2/2}$
shows that the operator-norm ball of radius $a/2$ has probability greater
than $1/2$. Its Frobenius distance from $W$ is
$D(W)=[\sum_j(s_j(W)-a/2)_+^2]^{1/2}$. Gaussian isoperimetry gives
$\mathbb P\{D(W)>r\}\le e^{-r^2/(2\sigma^2)}$.
Since $s\le2(s-a/2)$ when $s>a$, one has
$\|W_{>a}\|_F\le2D(W)$. Take $r=\sqrt K/2$.
\end{proof}

For $K\ge1$, define
\[
 a_K=\left[16L\lambda^L(K+1)(\sqrt K+1)^{L-2}\right]^{-1},
 \qquad J_K=(L+3)K.
\]
Use the auxiliary class $\mathcal F(J)$ from Corollary~\ref{cor:count};
it includes the Boolean functions admitting a strict margin-$1/2$,
depth-$L$, fixed-$\phi$ implementation of energy at most $J$ at some
finite width.

\begin{proposition}[Aggregate functional tail]
\label{prop:unbounded-aggregate}
Suppose $\gamma>1$ and
\[
 N^{\gamma-1}\ge8(d+1),\qquad
 N^{\gamma-1}\ge64a_K^{-2}.
\]
Then
\[
 \mu_{N,\gamma}\!\left(\bigcup_{g\notin\mathcal F(J_K)}U_{N,g}\right)
 \le L e^{-KN^\gamma/8}.
\]
\end{proposition}
\begin{proof}
Outside an event of probability at most $Le^{-KN^\gamma/8}$,
\[
 \frac dN\|W^1\|_F^2\le K,\quad N\|v\|_2^2\le K,
 \quad\|(W^s)_{>a_K}\|_F^2\le K\quad(2\le s<L).
\]
Indeed, the first and last random quantities are respectively
$N^{-\gamma}\chi^2_{N(d+1)}$ and $N^{-\gamma}\chi^2_N$.
For $u\ge8D$, Markov's inequality with exponential parameter $3/8$ gives
$\mathbb P\{\chi_D^2>u\}\le 2^D e^{-3u/8}\le e^{-u/4}$; the first width condition
therefore suffices. The second width condition permits
Lemma~\ref{lem:unbounded-spectral-tail} on each internal matrix.

Replace each internal matrix by its $>a_K$ part. Retained operator norms
are at most $\sqrt K$ and discarded operator norms at most $a_K\le1$.
For the original hidden states write
$M_s=N^{-1/2}\sup_x\|h_w^s(x)\|_2$. The zero-at-zero and Lipschitz
assumptions imply
\[
 M_s\le\sqrt{2K}\,\lambda^s(\sqrt K+1)^{s-1}.
\]
In particular individual neurons need not be bounded, but their normalized
Euclidean norm is controlled uniformly over inputs. If $e_s$ is the
normalized Euclidean difference between original and truncated hidden
states, then
\[
 e_1=0,\qquad e_s\le\lambda\sqrt K\,e_{s-1}+\lambda a_K M_{s-1}.
\]
Induction gives, for $s\ge2$,
\[
 e_s\le a_K(s-1)\sqrt{2K}\lambda^s(\sqrt K+1)^{s-2}.
\]
The output discrepancy is consequently at most
\[
 a_K(L-2)\sqrt2K\lambda^{L-1}(\sqrt K+1)^{L-3}<1/8
\]
when $L\ge3$. It is zero when $L=2$.
A globally margin-separated original draw therefore retains strict margin
$3/8$ after truncation. Doubling the readout gives strict margin $3/4$
and energy at most $K+(L-2)K+4K=J_K$. Its Boolean function belongs
to $\mathcal F(J_K)$, proving the claimed aggregate tail.
\end{proof}

The small spectral cutoff requires a sufficiently strong scale. At mean
field, the Gaussian bulk has operator norm of constant order; the next
sections control its functional effect through covariance reconstruction.

% Self-contained mean-field analytic ingredients, 22 October 2026.
\section{Gaussian bulk, tail atoms, and stable exponential tilts}
\label{sec:meanfield-preliminaries}
Let $Q$ be a positive semidefinite matrix on a finite input set of size
$m$, with $\max_xQ_{xx}\le M$ and $M\ge1$. Its reproducing-kernel
Hilbert space $\mathcal H_Q$ is its range, with norm
$\|u\|_Q^2=u^\top Q^+u$. Thus
$|u_x|\le\sqrt M\|u\|_Q$. All input-set dependence in the quantitative
bounds below is through $\log(2m)$.

\begin{lemma}[Gaussian projection onto a bounded bulk]
\label{lem:bulk-recession}
Let $Z_1,\ldots,Z_N$ be independent $N(0,Q)$ vectors. Put
$B=\sqrt{2M\log(4m)}$ and $A=4(K+\log2)$, where $K>0$.
There are measurable decompositions $Z_i=a_i+u_i$ with
$\|a_i\|_\infty\le B$ and $u_i\in\mathcal H_Q$ such that
\[
 \Prb\left\{N^{-1}\sum_i\|u_i\|_Q^2>A\right\}\le e^{-NK}.
\]
\end{lemma}
\begin{proof}
The closed convex set $\{a\in\mathcal H_Q:\|a\|_\infty\le B\}$
has Gaussian probability at least $1/2$ by a union bound over the $m$
coordinates. Let $a(z)$ be its metric projection in $\mathcal H_Q$,
and let $u(z)=z-a(z)$. Gaussian isoperimetry on the support of $Q$ gives
$\Prb\{\|u(Z)\|_Q>s\}\le e^{-s^2/2}$.
Thus $\E e^{\|u(Z)\|_Q^2/4}\le2$ by integration of the tail.
Independence and exponential Markov give
$\exp(N\log2-NA/4)=e^{-NK}$. The projection is continuous and hence
measurable. For $Q=0$ take $a=u=0$.
\end{proof}

\begin{lemma}[General Lipschitz version for auxiliary activations]
\label{lem:bulk-lipschitz-auxiliary}
Suppose instead only that $\phi(0)=0$ and $\phi$ is globally
$\lambda$-Lipschitz. Keep $B,A$ and the Gaussian projections $Z_i=a_i+u_i$
from Lemma~\ref{lem:bulk-recession}, and take $T\geq2B$.
Let $r:\R^m\to[0,1]$ equal one on $\|z\|_\infty\le T$, and define
$Q_{\mathrm{bulk}}=N^{-1}\sum_i r(Z_i)^2\phi(Z_i)\phi(Z_i)^\top$.
Then, on the projection-good event, the empirical activation covariance is within
$6\lambda^2 BMA/T$ of its same bounded bulk plus
\[
 R=\frac1N\sum_i(1-r(Z_i)^2)\phi(u_i)\phi(u_i)^\top.
\]
Writing $\delta_i=\|u_i\|_{\mathcal H_Q}$ and omitting zero terms, this is
\[
 R=\frac1N\sum_i(1-r(Z_i)^2)\delta_i^2
       \psi_i(v_i)\psi_i(v_i)^\top,\qquad
 v_i=u_i/\delta_i,\quad \psi_i(t)=\phi(\delta_i t)/\delta_i.
\]
Every $v_i$ has unit RKHS norm, every $\psi_i$ is $\lambda$-Lipschitz
and zero at zero, and the nonnegative mixture weights sum to at most $A$.
The same good event bounds the next covariance diagonal by
$2\lambda^2 B^2+2\lambda^2 MA$.
\end{lemma}
\begin{proof}
Lipschitz continuity gives
$\|\phi(Z_i)-\phi(u_i)\|_\infty\leq\lambda B$ and
$\|\phi(u_i)\|_\infty\leq\lambda\sqrt M\delta_i$.
If $w_i=1-r(Z_i)^2>0$, then
$\sqrt M\delta_i\ge T-B\ge T/2$. Hence
\[
 N^{-1}\sum_iw_i\le4MA/T^2,\qquad
 N^{-1}\sum_iw_i\delta_i\le2\sqrt M A/T.
\]
The outer-product difference is therefore at most
$4\lambda^2 BMA/T+4\lambda^2 B^2MA/T^2\leq6\lambda^2 BMA/T$.
The normalization is an exact algebraic identity and preserves the
Lipschitz constant. The diagonal bound follows from
$\|\phi(Z_i)\|_\infty\leq\lambda B+\lambda\sqrt M\delta_i$.
\end{proof}

\paragraph{A smooth global cutoff.}
Fix once and for all a smooth nonincreasing function
$\chi:\mathbb R\to[0,1]$ equal to one on $(-\infty,2]$ and zero on
$[5,\infty)$. For $T\geq1$, set
\[
 a=\frac{\log(2m)}{T^2},\quad
 s(z)=a^{-1}\log\sum_{j=1}^m e^{az_j^2},\quad
 r_T(z)=\chi(s(z)/T^2).
\]
Since $\|z\|_\infty^2\leq s(z)\leq\|z\|_\infty^2+T^2$,
we have $r_T=1$ on $\|z\|_\infty\leq T$ and $r_T=0$ on
$\|z\|_\infty\geq\sqrt5T$. Direct differentiation gives absolute constants
depending only on $\chi$ such that
\[
 \sup_z\sum_j|\partial_jr_T(z)|\leq C/T,
 \qquad
 \sup_z\sum_{j,k}|\partial_{jk}r_T(z)|
       \leq C\log(2m)/T^2.
\]
Indeed, on the support of the derivatives of $\chi$, every $|z_j|$ is at
most $\sqrt5T$. If $p_j=e^{az_j^2}/\sum_ke^{az_k^2}$, then
$\partial_js=2p_jz_j$ and
$\partial_{jk}s=2p_j\mathbf1_{j=k}
 +4a(p_jz_j^2\mathbf1_{j=k}-p_jp_kz_jz_k)$.
Their entrywise derivative sums are at most $2\sqrt5T$ and
$2+40\log(2m)$, respectively. The chain rule gives the assertion.

For any globally $\lambda$-Lipschitz activation $\phi$ with $\phi(0)=0$,
choose a fixed smooth nonnegative probability density $\rho$ supported
on $[-1,1]$ and mollify
$\psi_\tau(t)=\int\phi(t-\tau y)\rho(y)\,dy$, where $0<\tau\leq1$.
Then $\|\psi_\tau-\phi\|_\infty\leq\lambda\tau$,
$|\psi'_\tau|\leq\lambda$, and
$|\psi''_\tau|\leq\lambda\|\rho'\|_1/\tau$.
For ReLU one may alternatively use
$(t+\sqrt{t^2+\tau^2})/2$, with uniform error at most $\tau/2$.
For GELU one may use $\psi=\phi$ directly, since its first two
derivatives are bounded by absolute constants. Write
$b_x(z)=r_T(z)\psi(z_x)$.
These bounded features have uniform derivative bounds
\begin{align*}
 B_0&:=\max_x\|b_x\|_\infty\leq C\lambda T,\\
 B_1&:=\max_x\sup_z\sum_j|\partial_jb_x(z)|\leq C\lambda,\\
 B_2&:=\max_x\sup_z\sum_{j,k}|\partial_{jk}b_x(z)|
       \leq C\lambda(\log(2m)/T+\tau^{-1}).
\end{align*}
where for GELU one can set $\tau=1$ in the last bound. In general,
the empirical covariance of $b(Z_i)$ differs from $Q_{\mathrm{bulk}}$
by at most $C\lambda^2 T\tau$ in maximum norm, deterministically.


\begin{lemma}[A sparse exponential-tilt certificate]
\label{lem:sparse-tilt}
Let $F=(F_1,\ldots,F_D)$ be bounded observables of a probability law
$\mu$, with $|F_j|\le F_0$, and let $\widehat F$ be their empirical
means on $N$ independent draws. For $G,K>0$ and a positive integer $s$,
except on an event of probability at most
\[
 \exp\{s\log(2D+1)-N(K-G^2F_0^2/(2s))\},
\]
there is a tilt $d\mu_\Gamma=e^{\Gamma\cdot F-\Lambda(\Gamma)}d\mu$
with $\|\Gamma\|_1\le G$, relative entropy at most $K$, and
$\|\widehat F-\E_{\mu_\Gamma}F\|_\infty\le K/G$.
Here $\Lambda(\Gamma)=\log\E_\mu e^{\Gamma\cdot F}$.
\end{lemma}
\begin{proof}
Write $g(\Gamma)=\Gamma\cdot\widehat F-\Lambda(\Gamma)$.
For every fixed $\Gamma$, independence gives $\E e^{Ng(\Gamma)}=1$.
Every vector in the radius-$G$ $\ell^1$ ball is the mean of a random
vertex from $\{0,\pm Ge_1,\ldots,\pm Ge_D\}$. Average $s$ independent
such vertices to obtain $\Gamma_s$ with $\E\Gamma_s=\Gamma$.
For every fixed observation, Hoeffding's exponential lemma gives
\[
 \E_{\Gamma_s}e^{(\Gamma_s-\Gamma)\cdot F}
 \le e^{G^2F_0^2/(2s)}.
\]
Jensen's inequality and Fubini therefore show
$\E_{\Gamma_s}\Lambda(\Gamma_s)\le\Lambda(\Gamma)+G^2F_0^2/(2s)$.
For at least one of the at most $(2D+1)^s$ sparse averages,
$g(\Gamma_s)\ge g(\Gamma)-G^2F_0^2/(2s)$.
Exponential Markov and a union bound prove that
$\max_{\|\Gamma\|_1\le G}g(\Gamma)\le K$ outside the stated event.

Choose a maximizer $\Gamma_*$ on the compact ball. The first-order
optimality condition says
$\Gamma_*\cdot\nabla g(\Gamma_*)=G\|\nabla g(\Gamma_*)\|_\infty$,
including an interior maximum, where both sides are zero.
Concavity and $g(0)=0$ give
$g(\Gamma_*)\ge\Gamma_*\cdot\nabla g(\Gamma_*)$.
Since $\nabla g=\widehat F-\E_{\mu_\Gamma}F$, this proves the required
approximation. Finally the tilt entropy is
\[
 D(\mu_{\Gamma_*}\|\mu)
 =g(\Gamma_*)-\Gamma_*\cdot\nabla g(\Gamma_*)\le K.
\]
\end{proof}

\begin{lemma}[Polynomial stability of the tilted bounded bulk]
\label{lem:bulk-tilt-stability}
Let $b_x$ be any bounded $C^2$ features satisfying the preceding
$B_0,B_1,B_2$ bounds. For a matrix $\Gamma$ with
$\|\Gamma\|_{1,\mathrm{ent}}\leq G$, define
\[
 \mathcal F_\Gamma(Q)=
 \frac{\mathbb E_{Z\sim\mathcal N(0,Q)}
       [b(Z)b(Z)^\top e^{b(Z)^\top\Gamma b(Z)}]}
      {\mathbb E_{Z\sim\mathcal N(0,Q)}
                         e^{b(Z)^\top\Gamma b(Z)}}.
\]
Put $F_0=B_0^2$, $F_1=2B_0B_1$, and
$F_2=2B_1^2+2B_0B_2$. For all positive semidefinite $Q,Q'$,
\[
 \|\mathcal F_\Gamma(Q)-\mathcal F_\Gamma(Q')\|_{\max}
 \leq C_G\|Q-Q'\|_{\max},
 \quad
 C_G=\tfrac12\{F_2+2GF_1^2+2F_0(GF_2+G^2F_1^2)\}.
\]
Moreover, with $s$ any positive integer, except on an event of probability
\[
 \exp\{s\log(2m^2+1)-N(K-G^2F_0^2/(2s))\},
\]
the empirical covariance of $b(Z_1),\ldots,b(Z_N)$ is within $K/G$
in maximum norm of some $\mathcal F_\Gamma(Q)$ with
$\|\Gamma\|_{1,\mathrm{ent}}\leq G$ and tilt entropy at most $K$.
\end{lemma}
\begin{proof}
Apply Lemma~\ref{lem:sparse-tilt} to the $D=m^2$ bounded observables
$F_{xy}(z)=b_x(z)b_y(z)$. This gives the approximation and exceptional
probability, including the entropy bound.

For stability, interpolate $Q_t=(1-t)Q+tQ'$, set
$V=b^\top\Gamma b$, and let $\mathbb E_t$ denote Gaussian expectation
tilted by $e^V$. Gaussian differentiation and the quotient rule give
\begin{align*}
 \frac d{dt}\mathbb E_tF
 =\frac12\sum_{j,k}(Q'-Q)_{jk}\bigl\{
 &\mathbb E_t[\partial_{jk}F+
             \partial_jF\partial_kV+\partial_kF\partial_jV]\\
 &+\operatorname{Cov}_t(F,\partial_{jk}V+
                                   \partial_jV\partial_kV)\bigr\}.
\end{align*}
For $F=b_xb_y$, its sup norm and first/second derivative entry sums are
at most $F_0,F_1,F_2$. The corresponding derivative sums for $V$ are
at most $GF_1,GF_2$. Taking absolute values and summing proves the
claimed bound. All functions and derivatives are bounded; adding
$\epsilon I$ and passing to the limit proves the statement for singular
covariances as well. No reciprocal partition-function estimate is used,
so the stability constant is polynomial, not exponential, in $G$.
\end{proof}

\subsection{Stable finite realization of tail atoms}
The next certificate avoids any lower bound on the nonzero eigenvalues
of the covariance. It is used for both tail atoms and the final readout.

\begin{lemma}[An $\ell^1$ certificate for an RKHS vector]
\label{lem:rkhs-atom-certificate}
If $u\in\mathcal H_Q$, $\|u\|_Q\le1$, and $\eta>0$, there is a vector
$\beta\in\mathbb R^{\mathcal X}$ such that
\[
 \|Q\beta-u\|_\infty\le\eta,\qquad
 \|Q\beta\|_Q\le1,\qquad \|\beta\|_1\le\frac1{4\eta}.
\]
\end{lemma}
\begin{proof}
In the finite-dimensional Hilbert space $\mathcal H_Q$, minimize
$\tfrac12\|v\|_Q^2$ subject to $\|v-u\|_\infty\le\eta$.
A minimum exists, and $u$ is a strictly feasible point. The multiplier
conditions for the $2|\mathcal X|$ linear inequalities therefore apply.
With the signs chosen so that $v=Q\beta$, complementary slackness says
\[
 \beta\cdot(u-v)=\eta\|\beta\|_1.
\]
The left side is $\langle v,u\rangle_Q-\|v\|_Q^2$, and hence is at
most $\|v\|_Q-\|v\|_Q^2\le1/4$. The other two claims follow from
feasibility and the comparison with the feasible point $u$.
The argument also covers singular $Q$, since the optimization takes place
in $\mathcal H_Q$ itself.
\end{proof}

\begin{proposition}[Stable realization with auxiliary activations]
\label{prop:generic-atom-realization}
Let $Q$ have diagonal entries at most $M\ge1$, let $A\ge1$ and
$0<\varepsilon\le1$, and suppose
\[
 R=\sum_{j=1}^r c_j\psi_j(v_j)\psi_j(v_j)^\top,\quad
 c_j\ge0,\quad\sum_jc_j\le A,\quad\|v_j\|_Q\le1,
 \quad\psi_j(0)=0,\quad\operatorname{Lip}(\psi_j)\le\lambda.
\]
Take
\[
 \eta=\frac{\varepsilon}{64A\lambda^2(\sqrt M+1)},\qquad
 n=\left\lceil128A^2\lambda^4M^2\varepsilon^{-2}\log(4m^2)\right\rceil.
\]
Then $R$ admits an approximation $R'(Q')$ within $\varepsilon$ in maximum
norm, uniformly over all positive semidefinite $Q'$ satisfying
$\|Q'-Q\|_{\max}\le\eta^2$, by $n$ auxiliary neurons of the stated
activation class. The covariance coefficients are $A/n$, and the
incoming weights have the explicit bounds established below.
\end{proposition}
\begin{proof}
Sample pairs $(v_j,\psi_j)$ with probabilities $c_j/A$, adding a zero
pair for the remaining probability. The absolute entries of every
sampled outer product are at most $\lambda^2M$, so Hoeffding's inequality
gives samples whose average times $A$ approximates $R$ within
$\varepsilon/4$ entrywise. Apply the RKHS certificate to the sampled
$v_j$, obtaining $\|\beta_j\|_1\le1/(4\eta)$ and
$\|Q\beta_j-v_j\|_\infty\le\eta$.
For every $Q'$ with $\|Q'-Q\|_{\max}\le\eta^2$,
\[
 \|Q'\beta_j-v_j\|_\infty\le5\eta/4,
 \qquad
 \|\psi_j(Q'\beta_j)-\psi_j(v_j)\|_\infty\le5\lambda\eta/4.
\]
It follows that
\[
 R'(Q')=\frac A n\sum_j
             \psi_j(Q'\beta_j)\psi_j(Q'\beta_j)^\top
 \quad\hbox{satisfies}\quad
 \|R'(Q')-R\|_{\max}\le\varepsilon.
\]
The displayed constants leave ample room for the outer-product error.

If $Q'=\sum_\ell a_\ell h_\ell h_\ell^\top$ is a finite weighted
feature realization, each new neuron has activation $\psi_j$ and weights
\[
 W_{j\ell}=a_\ell\sum_y\beta_j(y)h_\ell(y),
 \qquad |W_{j\ell}|\le a_\ell\|h_\ell\|_\infty/(4\eta).
\]
Its covariance coefficient is $A/n$.
\end{proof}
These are ordinary finite
feedforward computations with neuron-specific activations; the common
capacity theorem controls them uniformly without coding the activations.
All intermediate values and coefficients obey polynomial bounds whenever
the preceding realization does.

% Deep mean-field closure, 22 October 2026.
% Input after the bulk-tail, RKHS-certificate, and generic-atom fragments.
\section{The deep mean-field endpoint}
\label{sec:meanfield-generic-closure}
Throughout this section $\phi(0)=0$ and $\phi$ is $\lambda$-Lipschitz,
where $\lambda\ge1$ is fixed. The Bayesian model uses this same
activation at every neuron. The finite reconstructions used only for
counting may have different zero-at-zero $\lambda$-Lipschitz activations
at different neurons; the preceding capacity theorem covers precisely
this enlarged class. No claim that those auxiliary activations are
additional parameters of the Bayesian learner is made.

We now assemble the local covariance estimates into a finite network,
working backward to choose error tolerances and forward to reconstruct
the layers. At $\gamma=1$, let $H=L-1$, $m=2^d$, and write
\[
 Q_0(x,y)=d^{-1}\bar x\cdot\bar y,\qquad
 Q_s(x,y)=N^{-1}\sum_{i=1}^Nh_i^s(x)h_i^s(y).
\]
Conditionally on the preceding layers, the $N$ preactivation fields at
layer $s$ are independent Gaussians with covariance $Q_{s-1}$. This is
exact at every finite width.

\subsection{One-step reconstruction and its polynomial modulus}
Here a modulus specifies how accurately the incoming covariance must be
reconstructed to achieve a prescribed outgoing accuracy. Its reciprocal
will be polynomially bounded, allowing iteration at fixed depth without
an exponential dependence on $m$. Fix $K\ge1$, put
\[
 q=4K,\qquad A=4(q+\log2),
\]
and suppose the input covariance has diagonal at most $M\ge1$.
For a desired output accuracy $0<\alpha\le1$, set
\[
 B=\sqrt{2M\log(4m)},\qquad
 T=\max\{1,2B,96\lambda^2BMA/\alpha\}.
\]
Use the smooth cutoff from the preceding section. Mollify $\phi$ at
scale $\tau\le1$ small enough that the maximum-norm difference between
its cutoff covariance and that using the actual activation is at most
$\alpha/16$, uniformly in the Gaussian covariance. For example
$\tau=\min\{1,\alpha/[64C\lambda^2(T+1)]\}$ works with the absolute
constant $C$ in the smoothing estimate. Let $F_0,F_1,F_2$ be the bounded
observable and derivative budgets in that section, and put
\[
 G=16q/\alpha,\qquad
 s=\max\{1,\lceil G^2F_0^2\rceil\},
\]
\[
 C_G=\tfrac12\{F_2+2GF_1^2+2F_0(GF_2+G^2F_1^2)\},\qquad
 \eta=\frac{\alpha}{512A\lambda^2(\sqrt M+1)},
\]
\[
 r(\alpha)=\min\{1,\eta^2,\alpha/[16(1+C_G)]\},\qquad
 N_0(\alpha)=\left\lceil\frac{2s\log(2m^2+1)}q\right\rceil.
\]
Every displayed quantity is bounded by a polynomial in
$K,M,\log(2m),\alpha^{-1}$, or, for $r$, its reciprocal is so bounded.
The constants may depend on $\lambda$ and the fixed smooth cutoff.

For $N\ge N_0(\alpha)$, except with probability at most $2e^{-NK}$,
the empirical outgoing covariance is within $3\alpha/16$ of
\[
 \mathcal F_\Gamma(Q)+R,\qquad
 \|\Gamma\|_{1,\mathrm{ent}}\le G,
\]
where $R$ is a mixture of unit-RKHS auxiliary-activation atoms with total
coefficient mass at most $A$. Indeed, the tail decomposition error is at
most $\alpha/16$, smoothing costs at most $\alpha/16$, and the entropy
certificate costs at most $q/G=\alpha/16$. The projection failure is
at most $e^{-Nq}$; the entropy failure is at most
\[
 \exp\{s\log(2m^2+1)-N(q-1/2)\}\le e^{-NK}.
\]
These bounds hold conditionally and uniformly in the input covariance.

Now suppose $Q'$ is a positive semidefinite covariance with
$\|Q'-Q\|_{\max}\le r(\alpha)$, represented by finitely many weighted
features from the preceding layer. The selected $\Gamma$ and atom
certificates remain fixed while $Q$ is replaced by $Q'$. The bulk change
is at most $\alpha/16$. The generic atom construction, with target
accuracy $\alpha/8$, uses
\[
 n_A=\left\lceil8192A^2\lambda^4M^2\alpha^{-2}\log(4m^2)\right\rceil
\]
auxiliary neurons and approximates $R$ within $\alpha/8$ at covariance
$Q'$. Its certificate parameter is exactly the $\eta$ above.

To realize the bulk, sample independent Gaussian rows of covariance
$Q'$ tilted by the bounded smooth observable. Evaluate the actual
activation $\phi$ at those rows and absorb the scalar cutoff into their
covariance coefficients. Each retained cutoff feature is bounded by
$B_b=\lambda\sqrt5T$. Take
\[
 n_B=\left\lceil1024B_b^4\alpha^{-2}\log(16m^2)\right\rceil.
\]
Hoeffding's inequality and a union bound over $m^2$ entries give a
realization with empirical covariance within $\alpha/16$ of its own
expectation. The expectation differs from
$\mathcal F_\Gamma(Q')$ by at most $\alpha/16$ because of the smooth
activation approximation. Thus all the accumulated errors are less than
$\alpha$. This construction uses $n_A+n_B$ neurons, a polynomial number
in $K,M,\log(2m),\alpha^{-1}$.

\paragraph{Finite weights and genuine feedforward evaluation.}
Write the preceding covariance as
$Q'=\sum_{\ell=1}^r a_\ell h_\ell h_\ell^\top$, with
$a_\ell\le V:=\max\{1,A\}$ and $\|h_\ell\|_\infty\le B_h$.
A bulk preactivation is realized with weights
$\sqrt{a_\ell}\xi_\ell$, where the Gaussian vector $\xi$ is tilted by
a density ratio at most $e^{2GF_0}$. Intersecting the empirical covariance
concentration event with a Gaussian-coordinate tail event ensures
\[
 |\xi_\ell|\le R_\xi
 :=\sqrt{4GF_0+2\log(16n_B\max\{r,1\})}
\]
for every sampled coordinate. The latter event fails with probability
at most $1/8$, whereas the preceding covariance event fails with
probability less than $1/8$, so the intersection has positive probability.
Bulk weights are therefore bounded by $\sqrt V R_\xi$.
The atom weights are $a_\ell\sum_y\beta_j(y)h_\ell(y)$ and are bounded
by $VB_h/(4\eta)$. All new covariance coefficients are at most $V$:
they are $r_T(z)^2/n_B$ for bulk rows and $A/n_A$ for atom rows.

For example, if $B_W$ denotes the maximum of the two incoming-weight
bounds, the actual hidden coordinates satisfy the deterministic bound
$B_h^{\rm new}\le\lambda rB_WB_h$. This bounds the untruncated physical
activations as well as the weighted features. Hence at any fixed number
of layers both the weights and all physical hidden coordinates remain
polynomially bounded. The cutoff scalars, tilt matrices, and RKHS
certificates select finitely many real weights; they are not evaluated
on the input when the resulting network predicts.
At the first layer use input features $h_\ell(x)=\bar x_\ell$ and
coefficients $a_\ell=1/d$. In particular an atom row is the ordinary
linear form with weights $d^{-1}\sum_y\beta_j(y)\bar y$.

\subsection{A uniform functional exceptional-mass bound}
\begin{proposition}[Deep mean-field functional compression]
\label{prop:generic-meanfield-compression}
For each fixed $L$ and $\lambda$, there are polynomials
$P_{L,\lambda}$ and $J_{L,\lambda}$ such that the following holds for
every $d\ge1$, $K\ge1$, and
$N\ge P_{L,\lambda}(d,K)$. Except on an event of raw Gaussian probability
at most $(2H+1)e^{-NK}$, every globally strict-margin draw has the same
Boolean sign function as an auxiliary depth-$L$ network with normalized
energy at most $J_{L,\lambda}(d,K)$. All its neuronwise activations vanish
at zero and have Lipschitz constant at most $\lambda$.
\end{proposition}
\begin{proof}
Define layerwise deterministic diagonal bounds
\[
 M_s=2[8\lambda^2\{\log(4m)+A+1\}]^s,\qquad 0\le s\le H.
\]
At layer $s$, apply the one-step construction with $M=M_{s-1}$ and
its corresponding projection cutoff $B_s$. On its projection-good event,
the next covariance has diagonal at most
$2\lambda^2B_s^2+2\lambda^2M_{s-1}A\le M_s$. Starting from
$\max_xQ_0(x,x)\le2$, induction therefore validates the diagonal budget
used at each layer. Every $M_s$ is polynomial in $d,K$ for fixed
$L,\lambda$ because $\log m=d\log2$. There are at most two
exceptional events per hidden layer, as accounted for in the one-step
construction.

Set the desired final covariance error to $\zeta_H=1/(128K)$.
Working backwards, for $s=H,H-1,\ldots,1$, use the one-step construction
with $\alpha=\zeta_s$ and $M=M_{s-1}$, and define
$\zeta_{s-1}=r_s(\zeta_s)$ using that layer's modulus.
Each $\zeta_s^{-1}$ is polynomially bounded in $d,K$ at fixed $L$:
this is a fixed finite composition of the polynomial reciprocal moduli
above. Taking the maximum of the finitely many thresholds
$N_{0,s}(\zeta_s)$ defines a polynomial width threshold
$P_{L,\lambda}(d,K)$ (increased to at least one if necessary).
Starting from the exact $Q_0$, reconstruct the covariances forwards.
The transported certificates are fixed for the original covariances;
their estimates hold uniformly for every admissible preceding
reconstruction. Thus no independence assumption is imposed on those
certificates. Induction yields a weighted-feature covariance
$\widehat Q_H$ with
\[
 \|\widehat Q_H-Q_H\|_{\max}\le1/(128K).
\]
The construction above bounds every layer width, every physical weight,
and every physical activation by polynomials in $d,K$ at fixed depth.

For the actual Gaussian readout $v_i\sim N(0,N^{-2})$, put
$E_v=N\|v\|_2^2=N^{-1}\chi_N^2$. Markov's inequality at exponential
parameter $1/4$ gives
\[
 \Pr\{E_v>8K\}\le\exp[-N(2K-\tfrac12\log2)]\le e^{-NK}.
\]
On its complement the original output $z$ belongs to $\mathcal H_{Q_H}$
and $\|z\|_{Q_H}^2\le8K$. The scaled form of
Lemma~\ref{lem:rkhs-atom-certificate}, with approximation accuracy $1/8$,
gives
\[
 \|Q_H\beta-z\|_\infty\le1/8,\qquad \|\beta\|_1\le16K.
\]
Replacing $Q_H$ by $\widehat Q_H$ adds at most $1/8$ to this error.
For $\widehat Q_H=\sum_i a_i h_ih_i^\top$, the readout
$v_i'=2a_i\sum_y\beta(y)h_i(y)$ therefore represents a function with
the same signs as $z$ and strict margin greater than $1/2$.
Its weights are bounded by $32KVB_h$.

Pad the finitely many reconstructed layers to their maximum width $r_*$
with zero neurons, which is harmless since every activation vanishes at
zero. If $B_*$ bounds all weights, the normalized energy is at most
\[
 [d(d+1)+(L-1)r_*^2]B_*^2.
\]
Taking a polynomial upper bound on this expression defines
$J_{L,\lambda}(d,K)$. A conditional union bound over the $H$ layers and
the readout proves the stated raw exceptional probability. The auxiliary
capacity theorem applies uniformly to the possibly different activations;
no finite coding of activation formulas is needed.
\end{proof}

% Uniform all-gamma corollary, 22 October 2026.
\section{One polynomial for every scaling at least as strong as mean field}
\label{sec:uniform-supercritical}
The endpoint and spectral estimates control complementary scale ranges.
For the remaining bounded widths, a total-energy bound suffices. Placing
all three ranges in one auxiliary function class produces one sample
event for every width and scale. Parameterize the raw
Gaussian prior by $0<t\le1$, with coordinate variances
\[
 \frac{t^2}{d}\quad\hbox{in the first layer},\qquad
 \frac{t^2}{N}\quad\hbox{in internal layers},\qquad
 \frac{t^2}{N^2}\quad\hbox{at the readout}.
\]
The scaling exponent $\gamma\ge1$ corresponds to
$t=N^{(1-\gamma)/2}$. Allowing every $t\in(0,1]$ also permits the scale
to be chosen independently at each width. Write
$\nu_{N,t}=\mu_{N,\sigma=t/\sqrt N}$ for this raw law.

\begin{lemma}[Uniformity of the compression ingredients]
\label{lem:uniform-scale-ingredients}
Let $A_f=A$ be the constant in
Proposition~\ref{prop:unbounded-target-mass}. At every $N\ge C$ and
$0<t\le1$,
\[
 \nu_{N,t}(U_{N,f})\ge e^{-A_fN/(2t^2)}.
\]
The mean-field functional compression proposition remains true with the
same polynomials $P_{L,\lambda},J_{L,\lambda}$ and raw exceptional bound
$(2H+1)e^{-NK}$ whenever $N\ge P_{L,\lambda}(d,K)$, uniformly for
every $t\in(0,1]$.
Also, whenever
\[
 t\le a_K/8,\qquad t^{-2}\ge8(d+1),
\]
the spectral functional compression at energy $(L+3)K$ holds at every
width, with raw exceptional mass at most $Le^{-NK/(8t^2)}$.
\end{lemma}
\begin{proof}
The target-cylinder proof uses scalar Gaussian variance $t^2$ in its
constrained coordinates and $t\le1$ in its deterministic error bound;
replacing $N^\gamma$ by $N/t^2$ gives the first assertion directly.

In the endpoint argument the conditional preactivation covariance is
now $t^2Q_{s-1}$ and the first one is $t^2Q_0$. Apply the same one-step
construction to these covariances. Their diagonal budgets do not increase.
If $Q_{s-1}'$ approximates $Q_{s-1}$, then
\[
 \|t^2Q_{s-1}'-t^2Q_{s-1}\|_{\max}
 \le\|Q_{s-1}'-Q_{s-1}\|_{\max}.
\]
The weighted feature coefficients for the next preactivation are
multiplied by $t^2$. Thus Gaussian incoming weights are multiplied by
$t$, atom incoming weights by $t^2$, and all existing upper bounds on
weights and feature magnitudes remain valid. Initially the input-feature
coefficients are $t^2/d$. The readout energy is
$E_v=t^2\chi_N^2/N$, whose upper tail is no larger than at $t=1$.
Hence the same polynomial thresholds, compressed-energy bound, and
raw exceptional probability apply. Although the reconstructed networks
may depend on $t$, the single auxiliary energy class containing them does
not depend on it.

For spectral compression, the weighted first-layer energy and readout
energy are $(t^2/N)\chi^2_{N(d+1)}$ and $(t^2/N)\chi^2_N$;
the internal entries have variance $t^2/N$. The two stated inequalities
are exactly the hypotheses used in the spectral proof. Its singular-value,
RMS-activation, and readout estimates therefore give the displayed tail,
without a lower bound on $N$.
\end{proof}

Choose the following budgets, all independent of learner width and scale:
\begin{align*}
 K_s&=4A_f+8\log(2L/\varepsilon),&
 t_*&=\min\{a_{K_s}/8,[8(d+1)]^{-1/2}\},\\
 K_m&=2+A_f/t_*^2+\log\{2(2H+1)/\varepsilon\},&
 N_*&=\left\lceil\max\{1,P_{L,\lambda}(d,K_m)\}\right\rceil,\\
 D_*&=d+2+(L-2)N_*,&
 J_0&=2A_f+2D_*\log2+4\log(2/\varepsilon),\\
 J_*&=(L+3)K_s+J_{L,\lambda}(d,K_m)+J_0.
\end{align*}
Each budget, including $t_*^{-1}$, is polynomially bounded in
$d,C,\log(1/\varepsilon)$ for fixed $L,\lambda$.

\begin{theorem}[All widths and all scales at least as strong as mean field]
\label{thm:uniform-all-scales}
Fix a depth $L\ge2$ and an activation $\phi$ with $\phi(0)=0$ and
Lipschitz constant at most $\lambda$. Suppose a Boolean target has a
width-$C$, depth-$L$ implementation with weights bounded by $C$ and
strict global margin $1/2$. Let the original fixed-$\phi$ Gaussian prior
have the variances displayed above, condition it on an absolute margin
$1/2$ at every Boolean input, and then condition on exact training labels.
For every input distribution on the Boolean cube and every
$0<\varepsilon,\delta<1$, one event of training-sample probability at
least $1-\delta$ gives posterior average sign error at most $\varepsilon$
\emph{simultaneously for every integer $N\ge C$ and every $0<t\le1$},
provided
\[
 n\ge\frac2\varepsilon
 \left[\mathcal B(J_*)+\log\frac1\delta\right].
\]
In particular, one fixed-depth polynomial sample bound works for every
$\gamma\ge1$, including $\gamma=1$, without a polynomial degree or
constant that deteriorates as $\gamma\downarrow1$.
ReLU, exact GELU, SiLU, tanh, and every other fixed Lipschitz activation
vanishing at zero are covered whenever the stipulated teacher exists.
For the correct label to have posterior probability strictly greater
than $2/3$ outside a test set of mass at most $\varepsilon$, apply the
bound with average-error accuracy $\varepsilon/3$.
\end{theorem}
\begin{proof}
Let $\mathcal F_*$ be the single auxiliary strict-margin function class
of energy at most $J_*$. We show that at every $N\ge C$ and $0<t\le1$,
the raw probability of margin-separated functions outside this class,
divided by the raw target probability, is at most $\varepsilon/2$.

If $t\le t_*$, spectral compression gives the ratio bound
\[
 L\exp\left[-\frac N{t^2}(K_s/8-A_f/2)\right]
 \le\varepsilon/2.
\]
This applies at every width.
If $t\ge t_*$ and $N\ge N_*$, the uniform endpoint compression gives
\[
 (2H+1)\exp\left[-N\left(K_m-\frac{A_f}{2t^2}\right)\right]
 \le\varepsilon/2.
\]

For the remaining bounded widths $C\le N<N_*$, use total energy instead
of a union over the continuum of possible scales. The parameter count is
$p=N(d+2)+(L-2)N^2$, so $p/N\le D_*$. The normalized energy has law
$\kappa_N=(t^2/N)\chi_p^2$. Its exponential moment at $N/(4t^2)$ is
$2^{p/2}$, and therefore
\begin{align*}
 \frac{\nu_{N,t}\{\kappa_N>J_0\}}{\nu_{N,t}(U_{N,f})}
 &\le \exp\left\{-\frac N{t^2}(J_0/4-A_f/2)
                                      +\frac p2\log2\right\}\\
 &\le \exp\left\{-\frac N{t^2}
                  (J_0/4-A_f/2-D_*\log2/2)\right\}
 \le\varepsilon/2.
\end{align*}
Here $t\le1$ and $N\ge1$ justify the last two inequalities. A globally
margin-separated draw with $\kappa_N\le J_0$ itself belongs to the
enlarged auxiliary class. Thus all three regimes are contained in the
same $\mathcal F_*$ apart from the asserted relative mass.

The target must belong to $\mathcal F_*$, since otherwise the relative
mass would be at least one. Normalization by the global margin event
cancels, and conditioning on any noiseless training labels preserves the
upper bound $\varepsilon/2$ on posterior mass outside $\mathcal F_*$.
The capacity theorem bounds $\log|\mathcal F_*|$ by
$\mathcal B(J_*)$. A finite-class union bound ensures, on one sample
event of probability at least $1-\delta$, that every consistent member
has true error at most $\varepsilon/2$. That event refers only to the
single function class, so it works for the entire continuum of scales
and every learner width simultaneously. Adding the posterior mass of
the complement gives the claimed average-error bound.
\end{proof}


\section{Exact GELU implementations of Boolean circuits}\label{sec:gelu}
The preceding theorems apply to small implementing networks. To translate ordinary Boolean circuits into that premise, take a layered fan-in-two circuit with NOR, COPY, and NOT gates and one constant wire. Every layer has width at most $C$, counting that wire. Inputs are encoded as bits $b_i=(1+x_i)/2$.

For ReLU, a NOR gate is $\relu(1-a-b)$, a copy is $\relu(a)$, and NOT is $\relu(1-a)$ on Boolean inputs. These operations preserve exact bits; a constant wire of value one supplies offsets.

\begin{proposition}[Exact same-depth GELU compilation]\label{prop:gelu-circuit}
A layered Boolean circuit with $H$ stages and width $C$, counting its constant wire, has an exact GELU implementation with $H$ hidden layers, hidden width at most $2C$, and individual weights bounded by $4$. Its real-valued output can be chosen to equal the target sign $\pm1$ exactly. Consequently the learning theorems apply with teacher complexity at most $\max\{4,2C\}$.
\end{proposition}
\begin{proof}
Write $\tau=\Phi(1)$, $s=2\tau-1$, $A=\tau/s$, and $B=(1-\tau)/s$. The elementary bound $\tau>2/3$ gives $1<A<2$ and $0<B<1$. Direct substitution shows
\begin{equation}\label{eq:gelu-relu-grid}
 A\gelu(u)+B\gelu(-u)=\relu(u)\qquad(u\in\{-1,0,1\}).
\end{equation}
At $u=1$, the left side is $A\tau-B(1-\tau)=1$; at $u=-1$ it is $-A(1-\tau)+B\tau=0$; and at zero it vanishes.

Represent each gate by two GELU neurons with preactivations $u$ and $-u$. Decode the gate's bit as the linear combination with coefficients $A,B$. For NOR take $u=1-a-b$, for COPY $u=a$, and for NOT $u=1-a$. These always lie in $\{-1,0,1\}$. Substitute the preceding layer's linear decoders into each new preactivation; this implements all stages without an extra hidden layer. A single constant neuron has value $\tau$ and is propagated by incoming weight $1/\tau$, so it supplies the required constant one in each preactivation. The first stage uses the affine bit encoding above.

Two copies of the same gate input produce coefficients of size at most $2A<4$; all other internal coefficients have magnitude at most $2$ or $1/\tau<3/2$. The final linear readout $2b-1$ has coefficients $2A,2B,-1/\tau$, all bounded by $4$. The total width is at most twice the original width, and the output is exactly $\pm1$. Pad to width $\max\{4,2C\}$ when a teacher complexity convention also requires weights bounded by the width.
\end{proof}

The identity in \eqref{eq:gelu-relu-grid} is required only on the three gate preactivations, not on the whole real line. Thus this is an exact discrete compilation using fixed coefficients, with no large-gain approximation or homogeneity assumption. The derivative $\Phi(t)+t(2\pi)^{-1/2}e^{-t^2/2}$ has absolute value less than $2$, verifying the activation hypothesis directly.

\section{Scope and remaining questions}
The result separates two obstacles caused by unbounded activations. Large individual neuron values are handled by Euclidean norm control in the capacity argument. At intermediate exponents $1<\gamma<2$, large total Gaussian parameter energy is handled by retaining the part with large singular values; discarded matrices have small operator norm and their accumulated output error can be bounded at fixed depth.

At the deep mean-field endpoint, the Gaussian bulk has operator norm of constant order. The proof therefore keeps its covariance contribution, using a bounded tilt and separate tail atoms. The covariance reconstruction is polynomial at each fixed depth but repeatedly composes polynomial moduli. It does not give a polynomial-depth guarantee or a competitive numerical bound.

The surrogate networks used for counting may have different zero-at-zero Lipschitz activations at different neurons. This is a proof device, not a change to the Bayesian model. In particular, the argument does not establish a same-activation circuit-compression converse for the generic activation class.

The global-margin conditioning is also substantive. It converts the real-output capacity bound into a finite class of Boolean functions and lets us compare raw target mass with raw exceptional mass. Removing it, changing the likelihood, or learning real outputs requires additional arguments.

\paragraph{Provenance and verification.}
This paper is part of the collaborative Gaussian-prior learning program, including contributions from Claude. The capacity, spectral, mean-field reconstruction, and exact GELU compilation arguments have undergone independent AI-agent mathematical review. The main learning statements have Lean proofs; the repository's theorem-coverage record specifies the certified endpoints and their exclusions.

\begin{thebibliography}{9}
\bibitem{blumer} A. Blumer, A. Ehrenfeucht, D. Haussler, and M. K. Warmuth. Occam's razor. \emph{Information Processing Letters} 24(6):377--380, 1987.
\bibitem{golowich} N. Golowich, A. Rakhlin, and O. Shamir. Size-independent sample complexity of neural networks. \emph{Proceedings of COLT}, 2018.
\bibitem{borell} C. Borell. The Brunn--Minkowski inequality in Gauss space. \emph{Inventiones Mathematicae} 30:207--216, 1975.
\end{thebibliography}
\end{document}
